% Problem record op_aaf9791beced84e4. This TeX file is derived from
% database/problems_json/op_aaf9791beced84e4.json by scripts/sync-tex.mjs; edit the JSON record
% and rerun the script rather than editing this file.
\section{Exact spectral gap of random Pauli rotations}
\label{sec:op_aaf9791beced84e4}

\paragraph{Problem.}
Is the exact Haar-mixing spectral gap of uniformly random Pauli rotations equal to $2^n(2^n-3)/(8(4^n-1))$ for every $n\geq4$?

Let $n\geq4$ and let $\mathcal P_n:=\{I,X,Y,Z\}^{\otimes n}\setminus\{I^{\otimes n}\}$ be the nonidentity Hermitian Pauli strings. A random Pauli rotation chooses $P\in\mathcal P_n$ uniformly and $\theta\in[0,2\pi)$ uniformly, then applies $e^{i\theta P}$.

On $L^2(\mathrm{SU}(2^n),\mathrm{Haar})$, define the averaging operator and its spectral gap by

\begin{equation}
(K_nf)(U):=\mathbb E_{P,\theta}f(e^{i\theta P}U),
\qquad
\Delta_n:=1-\|K_n|_{L^2_0}\|_{2\to2},
\label{eq:aaf9-1}
\end{equation}

where $L^2_0$ is the subspace of functions with zero Haar mean. Each step is one Pauli rotation, not one nearest-neighbor two-qubit gate.

Is the exact gap given, for every $n\geq4$, by

\begin{equation}
\Delta_n=\frac{2^n(2^n-3)}{8(4^n-1)}?
\label{eq:aaf9-2}
\end{equation}

The displayed definitions, constraints, and target bounds are recorded in Eqs.~\eqref{eq:aaf9-1}, \eqref{eq:aaf9-2}.

\subsection*{Status}

Unsolved

\subsection*{Source}

The question is explicitly posed or retained as open in the cited primary literature \sourcecite{ref:aaf9-baer26}{Baer26}.  The statement is rewritten here to make its hypotheses and success criterion self-contained.

\subsection*{Progress}

\begin{itemize}
  \item Baer and Haah prove the bounds
  \begin{equation}
\frac{4^n+16}{16(4^n-1)}
  \leq\Delta_n
  \leq\frac{2^n(2^n-3)}{8(4^n-1)}.
\label{eq:aaf9-3}
\end{equation}
  The upper bound comes from the fourth-moment representation $U^{\otimes4}\otimes\overline U^{\otimes4}$; see Theorem 3.16 and Proposition 3.17. \sourcecite{ref:aaf9-baer26}{Baer26}

The displayed definitions, constraints, and target bounds are recorded in Eqs.~\eqref{eq:aaf9-3}.

  \item Proposition 3.17 settles $n\leq3$. In particular, $\Delta_3=5/63$, whereas the first unresolved instance has
  \begin{equation}
\frac1{15}\leq\Delta_4\leq\frac{26}{255}.
\label{eq:aaf9-4}
\end{equation}
  Conjecture 3.48 asserts that the upper bound is always attained. \sourcecite{ref:aaf9-baer26}{Baer26}

The displayed definitions, constraints, and target bounds are recorded in Eqs.~\eqref{eq:aaf9-4}.

  \item The 23 July 2026 preprint proves a constant gap, but explicitly leaves this exact formula conjectural. No subsequent resolution of Conjecture 3.48 was located. \sourcecite{ref:aaf9-baer26}{Baer26}
\end{itemize}

\subsection*{References}

\begin{enumerate}[label={},leftmargin=4.8em,labelsep=0.5em]
  \footnotesize
  \item[\textup{[Baer26]}]\label{ref:aaf9-baer26}
  T. Baer and J. Haah, "Random unitary circuits with constant spectral gap," arXiv preprint (2026), version 1, 23 July 2026. \href{https://arxiv.org/abs/2607.20919}{arXiv:2607.20919}.
\end{enumerate}

\subsection*{Comment}

The unresolved issue is whether any higher representation mixes more slowly than the identified fourth-moment mode. This would determine the exact spectral bottleneck across all moment orders, not merely establish convergence to Haar randomness.

\subsection*{Field}

Quantum algorithm

\subsection*{Topic}

Random circuit sampling; Quantum circuit complexity

\subsection*{ID}

\texttt{op\_aaf9791beced84e4}
